% Complete reciprocal finite proofs for ED26-BM-RAT2021-10-2.

\begin{annexitem}{annex:p17.rat102-pd}{thm:p17-rat102-pd}{Dominance, equilibrium and all four efficiency comparisons}
\begin{proof}
Against the column player's probability \(q\in[0,1]\) of cooperation,
the row player's pure expected utilities are
\[
                  U(c)=qR+(1-q)S,\qquad U(d)=qT+(1-q)P.
\]
Thus
\[
              U(d)-U(c)=q(T-R)+(1-q)(P-S)>0.
\]
Both differences are strictly positive; at the endpoints one coefficient
is one, and in the interior both coefficients are positive. The column
player has the symmetric strict comparison. If either player assigns
positive probability \(r\) to \(c\), that player's payoff is
\(rU(c)+(1-r)U(d)<U(d)\), so pure \(d\) is a profitable deviation.
Consequently any equilibrium must have both players choose pure \(d\).
At \((d,d)\), switching to \(c\) lowers payoff from \(P\) to \(S\),
as does any mixture with positive weight on \(c\). Hence \((d,d)\)
is an equilibrium and is the only independent mixed equilibrium.

Because \(R>P\), \((R,R)\) strictly improves both coordinates of
\((P,P)\), making \((d,d)\) inefficient. The pure outcome \((T,S)\)
is efficient because \(T\) is the first player's unique highest utility
among the four outcomes: no different outcome leaves that player at least
as well off. The symmetric argument applies to \((S,T)\). Finally
\((R,R)\) is efficient: \((P,P)\) lowers both utilities, \((T,S)\)
lowers the second, and \((S,T)\) lowers the first. All four pure
outcomes are accounted for.

No efficiency claim against every joint lottery has been made. Depending
on cardinal values, a correlated lottery over the asymmetric outcomes can
give both players more than \(R\). The result deliberately restricts
the efficiency domain to the four pure outcomes.
\end{proof}
\end{annexitem}

\begin{annexitem}{annex:p17.rat102-hilo}{thm:p17-rat102-hilo}{All Hi--Lo equilibria and the mixing calculation}
\begin{proof}
Against probability \(q\) of \(H\), the row player's payoffs are
\(U(H)=aq\) and \(U(L)=b(1-q)\). Equality is equivalent to
\(q=r=b/(a+b)\). Since \(a>b>0\), \(0<r<1/2\).
If \(q<r\), the unique best response is \(L\); if \(q>r\), it is
\(H\); at equality every mixture is optimal. The column player's
symmetric threshold concerns \(p=r\).

Against pure \(H\), \(H\) gives \(a>0\) rather than zero, so
\((1,1)\) is an equilibrium. Against pure \(L\), \(L\) gives
\(b>0\) rather than zero, so \((0,0)\) is an equilibrium. At
\((r,r)\), both support actions of each player give the same payoff,
so their displayed mixtures are optimal: it too is an equilibrium.

For exhaustiveness, let \((p,q)\) be an equilibrium. If \(p=0\),
the column player's unique best response is pure \(L\), forcing
\(q=0\). If \(p=1\), its unique response is pure \(H\), forcing
\(q=1\). Otherwise \(0<p<1\). A mixture putting positive weight
on both actions can be optimal only if their utilities agree: if one were
lower, removing its positive weight would strictly improve payoff.
Thus \(q=r\), which is interior, and the column player's nontrivial
mixture similarly forces \(p=r\). No other equilibrium exists.

At the interior equilibrium either pure action gives
\[
                        ar=b(1-r)=\frac{ab}{a+b}.
\]
Their mixture has the same expectation. Moreover,
\(1-r=a/(a+b)>r\) and \(ab/(a+b)=b\,a/(a+b)<b<a\).
The pure payoff \((a,a)\) improves both coordinates of \((b,b)\).
Since \(H\) is better against pure \(H\), while \(L\) is better
against pure \(L\), neither action strictly dominates the other.
All claims follow without assuming rational players select the superior
equilibrium or that actual players use these mixing probabilities.
\end{proof}
\end{annexitem}

\begin{annexitem}{annex:p17.rat102-pump}{thm:p17-rat102-pump}{Sequential cash loss and the acceptance boundary}
\begin{proof}
Let \(k\) be a nonnegative integer with \(w_0\geq3k\varepsilon\).
For \(k=0\), the result holds without trades. After \(j\) accepted
trades, induction gives cash \(w_j=w_0-j\varepsilon\): each step
subtracts the stated price. Before trade \(j+1\), for
\(0\leq j<3k\),
\[
             w_j\geq(3k-j)\varepsilon\geq\varepsilon.
\]
Each assumed acceptance therefore applies when needed. The alternative
sequence is \(x,z,y,x,z,y,x,\ldots\), so every three steps return
\(x\), and induction on rounds returns \(x\) after \(3k\) steps.
Cash is then \(w_0-3k\varepsilon\), a loss of \(3k\varepsilon\).

This argument depends on acceptance at the specified cash levels. An
unpriced cycle does not compare \((z,w-\varepsilon)\) with \((x,w)\).
For a concrete counterexample to that inference, extend the alternative
cycle to alternative--cash pairs by always preferring higher cash, and
using the cyclic alternative comparison only at equal cash. The equal-cash
comparison remains cyclic, but every proposed paid exchange is rejected
because it lowers cash. This relation is not asserted transitive: it only
shows that an ordinal cycle does not entail paid acceptance. The theorem
also makes no claim of unlimited repetition with finite cash, or that
real groups make these trades.
\end{proof}
\end{annexitem}

\begin{annexitem}{annex:p17.rat102-policy}{thm:p17-rat102-policy}{Budget feasibility, the votes and the priority procedure}
\begin{proof}
Dividing the no-deficit inequality by two gives \(d+o-1\leq t\).
If \(d=o=1\), the left side is one, and the binary \(t\) must
equal one. For the other three pairs \((d,o)\), the left side is
zero or minus one, so every binary \(t\geq0\) satisfies it.
This exhausts the binary pairs and proves equivalence to
\((d\land o)\Rightarrow t\).

For member 1, revenue change 2 finances spending change \(1+1=2\).
For member 2, zero revenue accompanies \(1-1=0\) spending change.
For member 3, zero revenue accompanies \(-1+1=0\). All three rows
satisfy the constraint. There is one vote for higher taxes and two
against, producing \(t=0\). Each spending increase has two votes
for and one against, producing \(d=o=1\). The majority row therefore
has spending change two against revenue zero and violates the constraint.

In the priority procedure, first-stage majority with zero on ties selects
binary \(d,o\) for any finite binary profile. If both are one,
\(t=d\land o=1\) is exactly the financing requirement. Otherwise
\(t=0\), which the complete preceding case analysis showed feasible.
Thus every first-stage pair yields a feasible final row. For the displayed
profile, both first-stage values are one and the final row is \((1,1,1)\).

The tax verdict is changed substantively: direct majority rejects the tax
increase, while the derived verdict accepts it because of the other
decisions. The procedure is not a common independent proposition-wise
rule. Feasibility does not establish fairness, truth, efficiency or
legitimacy. The budget equation is an editorial model of Schmid's
example, not fiscal data or an equation attributed to the publisher text.
\end{proof}
\end{annexitem}
